\documentclass[11pt, a4paper]{article}

\usepackage[a4paper, top=2.5cm, bottom=2.5cm, left=2cm, right=2cm]{geometry}
\usepackage{amsmath}
\usepackage{amssymb}
\usepackage[colorlinks=true, linkcolor=blue, urlcolor=blue, citecolor=blue]{hyperref}

\title{Theory of Everything - Volume 42 \\ \Large Stellar Engineering \& Kardashev Scales}
\author{Omega Protocol}
\date{\today}

\begin{document}

\maketitle

\section*{Layman's Summary}
How do you measure the power of an alien civilization? The Kardashev Scale categorizes them by energy use. Volume 42 proves that these scales aren't just science fiction categories; they are strict mathematical thresholds defined by the quantum limits of energy extraction.

\section{Kardashev Ascension (Axiom 45)}
We establish that as an informational network expands, its energy footprint naturally hits specific geometric limits: the planet, the star, and the galaxy. A Type I civilization controls a planet. Type II controls a star. Type III controls a galaxy. These are fundamental quanta of cosmic expansion.

\section{Dyson Swarm Dynamics (Theorem)}
To capture a star's energy, you don't build a solid metal sphere around it (which would be unstable). We prove that the mathematically optimal solution is a "Dyson Swarm"—millions of independent, unlinked orbital nodes buzzing around the star like bees. This maximizes energy capture while perfectly balancing gravity and light pressure.

\section{Star Lifting (Theorem)}
Stars naturally waste energy and burn out too fast. We prove that an advanced network can perform "Star Lifting"—using massive magnetic fields to physically pull excess gas off a star. This lowers the star's mass, slowing down its fusion rate, and extending its lifespan from billions to trillions of years.

\end{document}
